Zu Beginn der Projektarbeit haben wir, Wuyi und Arthur, uns zusammengesetzt und die Aufgaben, die uns als Gruppe \emph{Item- und Charakterdesign} zugeteilt wurden, klar definiert. Die Aufgaben sind folgende:

\begin{itemize}
	\item Festlegen der Anzahl an Lehrern und Items, sowie Auswählen der im Spiel vertretenen Lehrer für die jeweiligen Fachschaften
	\item Zeichnen der Spielerfigur sowie der Lehrer und im Spiel benutzten Items
	\item Implementieren des Spielers, der Lehrer und Items in das Programm 
\end{itemize}

\subsection{Spielerfigur}
Die Spielerfigur soll einerseits visuell den Benutzer im Spiel involvieren. Dafür gibt es über ein Menü die Möglichkeit, Haare, Hautfarbe und Kleidung individuell einzufärben.

Anderseits stellt sie den Hauptcharakter des Spieles dar und besitzt dementsprechend Funktionen, welche den Kampf gegen die Lehrer ermöglichen, wie auch ein Inventar, in welchem die erhaltenen Items gespeichert und mit Bild und Namen angezeigt werden. Dieses kann durch Anklicken der Spielerfigur geöffnet und mit einem Button wieder geschlossen werden. Die für den Kampf relevanten Funktionen wurden mit Absprache von der Gruppe \emph{Kampfkonzept} eingefügt und lagen außerhalb unserer Zuständigkeit.

\subsection{Lehrer*innen}
Die Lehrer und Lehrerinnen treten im Spiel als Antagonisten auf, die es zu besiegen gilt. Geplant waren 3 Lehrer*innen pro Fachschaft mit variierenden Schwierigkeitsgraden im Kampf. Um diese darzustellen, braucht es sowohl ein Bild des Lehrers/der Lehrerin, als auch Eigenschaften wie einen Namen, den Schwierigkeitsgrad oder das \glqq{}Level\grqq{}, die Zugehörigkeit zu einer Fachschaft, eine Menge an Lebenspunkten sowie eine Auswahl von Items, welche im Kampf zur Verfügung stehen.

\subsection{Items}
Die Items sind gebräuchliche Gegenstände/Materialien des Schulalltages oder auch bekannte Modelle/Konzepte der MINT-Fächer und erfüllen die Rolle von Waffen und Belohnungen im Spiel, die durch das erfolgreiche Beenden von Minispielen oder Besiegen von Lehrern nach einem Duell gewonnen und im Kampf gegen selbige auch eingesetzt werden können. Jedes Item besitzt einen Namen, ein Bild, welches im Spiel angezeigt wird, und einen Stärkegrad oder \glqq{}Level\grqq{}, das den Schaden im Kampf bestimmt. Es gibt 30 Items, sechs pro Fachschaft, mit jeweils zwei Items der Stärke null, zwei der Stärke eins und zwei der Stärke zwei. Nach dem Abschluss eines Minispiels erhält der Spieler ein Item der Fachschaft, zu dem das Spiel gehört. Der Stärkegrad des erhaltenen Items entspricht dabei dem Schwierigkeitsgrad jenes Spiels. Das heißt, je schwerer das Minispiel, desto höher ist das Level des erhaltenem Items.
